\documentclass[11pt]{article}
\usepackage[margin=1in]{geometry}
\usepackage{amsmath,amssymb,amsthm}
\usepackage[hidelinks]{hyperref}
\newtheorem{theorem}{Theorem}
\title{A square-clique lower bound for Closest String with a balanced alphabet}
\author{Atlas\thanks{The proof in this note was found using ChatGPT.}}
\date{}
\begin{document}
\maketitle

In \textsc{Closest String}, we are given explicitly listed strings of common length $L$ and a nonnegative integer radius $d$ in binary. The question is whether a string of length $L$ is within Hamming distance $d$ of every input string. Let $q$ be the number of distinct symbols used in the input strings, and let $N$ be the full binary input length. We use $[k]=\{1,\ldots,k\}$. The $k\times k$ \emph{Clique Hypothesis} \cite{szeider} excludes deterministic $2^{o(k\log_2 k)}$ time for deciding whether a graph on $[k]\times[k]$ contains a clique with one vertex per row.

\hypertarget{closest-string-balanced-kxk}{}
\begin{theorem}\label{thm:cs-kxk-balanced}
Under the $k\times k$ Clique Hypothesis, no deterministic algorithm solves \textup{\textsc{Closest String}} in time $2^{o(d\log_2 d)}N^{O(1)}$ on inputs satisfying $L=d+1$ and $q=2d+3$.
\end{theorem}

We use the reductions of Lokshtanov, Marx, and Saurabh~\cite{lms}. The promise on $q$ requires that each symbol actually appears. The sets requiring one selection per row in their intermediate instance ensure that all symbols appear, independent of the source graph's edges.

\begin{proof}
The deterministic construction in \cite[Corollary~2.6 and Lemma~2.8]{lms} maps a $k\times k$ \textsc{Clique} input to a disjunction of $2^{O(k\log\log k)}$ instances on $[k]\times[k]$ in total $2^{O(k\log\log k)}$ time. In those instances the clique must select one vertex from every row and column. Complement each graph to obtain the corresponding independent-set question. We may handle $k<3$ directly.

For one such graph, use two copies of its grid in separate diagonal blocks of a $2k\times 2k$ grid. Retain only cross-block conflicts. For each original row $i$ add a conflict between $(i,j)$ and $(i+k,j'+k)$ whenever $j\ne j'$. For each edge of the complemented graph with ordered endpoints $(i,j),(i',j')$, add a conflict between $(i,j)$ and $(i'+k,j'+k)$, using both orders. The copied selections in the two blocks must agree row by row because of the first conflicts, and the second conflicts forbid any edge within the original selection. Thus the original graph has the required independent set exactly when the two-block grid has a conflict-free selection confined to the diagonal blocks with one cell per row and column. This is the construction of \cite[Theorem~2.11]{lms}.

Let $h=2k$. Form a family of sets of cells in the new grid, following \cite[Theorem~2.12]{lms}. For each row $i\in[h]$, include the set of cells of that row in its diagonal block. For each conflict between two cells, include the union of these two row sets after removing the respective conflicting cell from each. Each set contains at most one cell per column. For the conflict sets this holds because the two rows have disjoint column ranges. Suppose a selection chooses one cell per column and hits every set. Hitting the $h$ disjoint row sets with $h$ selected cells forces one cell per row, all in diagonal blocks. A conflict set is then hit precisely when at least one of its two rows selects a cell different from the conflicting cell. Consequently the required selection exists if and only if a choice of one cell per column hits every set.

Convert each set $S$ to $h+1$ strings of length $h$, one for each dummy symbol $\delta\in\{h+1,\ldots,2h+1\}$, as in \cite[Theorem~3.1]{lms}. In position $j$ write the row index $i$ if $(i,j)\in S$, and write $\delta$ if no such cell belongs to $S$. Set the radius to $d=h-1$. A hitting selection gives a center by writing the chosen row at each column. This center agrees in at least one position with every input string. Conversely, a center of length $h$ omits some dummy symbol $\delta$. Its agreement with the string for $(S,\delta)$ identifies a selected cell of $S$ at a position where the center has an ordinary symbol. Assign any row to each position with a dummy symbol in the center. The resulting choice hits every set.

In the strings for row $i$ its own symbol $i$ appears in the $k$ columns of its block, and each dummy symbol appears in the strings for this row in the other $k$ columns. The used alphabet therefore has exactly $2h+1=2d+3$ symbols, and the common length is $h=d+1$. These promises hold for every query. There are polynomially many sets and strings per grid, and each symbol can be encoded using $O(\log(k+2))$ bits. Each transformation after the initial disjunction takes polynomial time in $k$. A hypothetical algorithm as in the theorem solves each query in $2^{o(k\log k)}k^{O(1)}$ time since $d=2k-1$. Summing over the disjunction still takes $2^{o(k\log k)}$ time, contrary to the hypothesis.
\end{proof}

\begin{thebibliography}{2}
\bibitem{lms}
Daniel Lokshtanov, D\'aniel Marx, and Saket Saurabh.
\emph{Slightly Superexponential Parameterized Problems}.
SIAM Journal on Computing 47(3), 675--702, 2018.
\href{https://doi.org/10.1137/16M1104834}{doi:10.1137/16M1104834}.
\bibitem{szeider}
Stefan Szeider.
\emph{After ETH: the $k\times k$ Clique Hypothesis}.
\href{https://doi.org/10.5281/zenodo.24088151}{Zenodo preprint v3}, 2026.
\end{thebibliography}
\end{document}
